\section{Тема 8. Геометрия и теория множеств}

\subsection{Постановка задачи (вариант~4)}

Из числового множества мощности $n$ выбрать все подмножества,
сумма элементов каждого из которых лежит в пределах от $a$ до $b$.

\subsection{Математическое обоснование}

Каждое подмножество $n$-элементного множества соответствует
$n$-разрядному двоичному числу (маске) от $0$ до $2^n-1$
(см. также Лаб.~2.1 по дискретной математике, метод двоичного
счёта). Для каждой маски вычисляется сумма выбранных элементов:
\[
  S(\text{mask}) = \sum_{i:\,\text{бит}_i(\text{mask})=1} x_i,
\]
и подмножество выводится, если $a \le S(\text{mask}) \le b$.
Сложность алгоритма~--- $O(n \cdot 2^n)$.

\subsection{Блок-схема алгоритма}

\begin{center}
\begin{tikzpicture}[
  node distance=10mm, start chain=going below,
  nodes={draw=black, top color=white, bottom color=lightgray!50,
         minimum width=5.6cm, align=center, font=\small},
  arrow/.style={->, >=Stealth, thick},
]
  \node[draw, rounded corners, on chain] {Start};
  \node[shape=trapezium, trapezium left angle=70,
        trapezium right angle=110, on chain, draw, join=by arrow]
       {$n$,\ $X[]$,\ $a$,\ $b$};
  \node[shape=chamfered rectangle, chamfered rectangle xsep=10pt,
        on chain, draw, join=by arrow]
       (loop) {$\mathrm{mask} := 0,\ 2^n{-}1$};
  \begin{scope}[local bounding box=scl]
    \node[shape=rectangle, on chain, draw, join=by arrow]
         {$S := \sum_{i:\,\mathrm{bit}_i} X[i]$};
    \node[shape=diamond, aspect=2, on chain, draw, join=by arrow,
          minimum width=5.2cm]
         (c) {$a \le S \le b$\,?};
    \node[shape=trapezium, trapezium left angle=70,
          trapezium right angle=110, draw, right=2cm of c]
         (out) {подмножество};
    \draw[arrow] (c.east) -- node[above,font=\scriptsize]{$+$} (out.west);
    \coordinate[on chain, join=by arrow] (le);
    \draw[arrow] (out.south) |- (le);
  \end{scope}
  \draw[arrow] (scl.south) -| ($(scl.west)-(0.9,0)$) |- (loop);
  \coordinate[on chain] (sp);
  \draw[arrow] (loop) -| ($(scl.east)+(0.9,0)$) |- (sp);
  \node[draw, rounded corners, on chain, join=by arrow] {End};
\end{tikzpicture}
\end{center}

\subsection{Текст программы}

\lstinputlisting[style=Cstyle]{../src/t08.c}

\subsection{Тестовые данные}

\begin{center}
\begin{tabular}{l|l}
\toprule
Множество, границы & Подмножества \\
\midrule
$\{1,2,3\}$, $a=2,b=4$ & $\{2\},\{3\},\{1,2\},\{1,3\}$ \\
\bottomrule
\end{tabular}
\end{center}

\textbf{Результат выполнения программы:}
\begin{lstlisting}[language={},basicstyle=\ttfamily\small,frame=single]
Введите мощность множества n: 3
Введите элементы множества:
1 2 3
Введите границы a b: 2 4
Подмножества с суммой в [2.0, 4.0]:
  {1,2}  сумма=3
  {3}  сумма=3
  {1,3}  сумма=4
  {2}  сумма=2
Найдено подмножеств: 4
\end{lstlisting}
